\section{Fabric Architecture}
\label{sec:arch}

The fabric is a $\text{ROWS}\times\text{COLS}=4\times4$ mesh of identical
processing elements operating on 16-bit signed words. All PEs share a single
clock domain and advance synchronously; a global \emph{step} enable, asserted
by the controller, gates every PE register so that the host advances the whole
array a deterministic number of cycles per command.

\subsection{Processing element}

Figure~\ref{fig:pe} shows the PE datapath. Each PE has two operand multiplexers
that independently select one of seven sources: the registered output of a
north, south, east or west neighbour; a 16-bit immediate carried in the
configuration word; the PE's own register (a feedback path); or a constant
zero. The selected operands drive an arithmetic-logic unit; its result is
latched into the single output register when \emph{step} is asserted.

\begin{figure}[t]
  \centering
  \resizebox{\columnwidth}{!}{% figures/pe.tex — processing-element datapath (logical).
\begin{tikzpicture}[node distance=6mm]
  % operand source list
  \node[align=right, font=\scriptsize, color=cWire] (srcs)
        {N, S, E, W\\ const (imm)\\ self ($\reg{r}$)\\ zero};

  % two operand muxes
  \node[mux, right=14mm of srcs, yshift=6mm]  (ma) {mux};
  \node[mux, right=14mm of srcs, yshift=-9mm] (mb) {mux};
  \node[lbl, above=0.5mm of ma] {sel\_a};
  \node[lbl, below=0.5mm of mb] {sel\_b};

  % alu and register
  \node[alu, right=12mm of $(ma)!0.5!(mb)$] (alu) {ALU};
  \node[reg, right=12mm of alu] (reg) {$\reg{r}$};
  \node[blk, below=10mm of alu, minimum width=30mm] (cfg)
        {config word\\{\scriptsize [imm | op | sel\_a | sel\_b]}};

  % wires: sources -> muxes
  \draw[wire] (srcs.east) -- (ma.west);
  \draw[wire] (srcs.east) -- (mb.west);
  % muxes -> alu
  \draw[wire] (ma.east) -- node[lbl,above]{a} ([yshift=3mm]alu.west);
  \draw[wire] (mb.east) -- node[lbl,below]{b} ([yshift=-3mm]alu.west);
  % alu -> reg -> out
  \draw[wire] (alu.east) -- (reg.west);
  \draw[wire] (reg.east) -- node[lbl,above]{dout} ++(12mm,0);
  % feedback self
  \draw[wire] (reg.south) |- ([yshift=-4mm]reg.south) -| (srcs.south)
        node[lbl, pos=0.25, below]{self};
  % config fan-out
  \draw[wire] (cfg.north) -- node[lbl,right]{op} (alu.south);
  \draw[wire] (cfg.west) -- ++(-6mm,0) |- (ma.south) node[lbl,pos=0.75,below]{};
  \draw[wire] (cfg.west) -- ++(-6mm,0) |- (mb.south);
  % step
  \draw[wire] ([xshift=-6mm]reg.south) ++(0,-8mm) -- node[lbl,left]{step} (reg.south);
\end{tikzpicture}
}
  \caption{Processing-element datapath: two operand muxes select from the four
  neighbours, the immediate, the self register or zero; the ALU result is
  latched by a single step-gated register.}
  \label{fig:pe}
\end{figure}

The ALU implements sixteen operations selected by a 4-bit opcode: a no-op that
holds the register; a pass-through; addition, subtraction and multiplication
(the low 16 bits of the product); a multiply-accumulate and a simple
accumulate, both folding into the register; the bitwise \textsc{and}/\textsc{or}
/\textsc{xor}; logical/arithmetic shifts; signed maximum, minimum and absolute
value; and a constant load. Multiplication and multiply-accumulate keep the low
half of the result, so arithmetic wraps modulo $2^{16}$, consistent with a
16-bit datapath.

A single 32-bit configuration word fully specifies a PE: the low sixteen bits
are the immediate, followed by the 4-bit opcode and the two 3-bit operand
selects. One such word per PE constitutes a complete array configuration.

\subsection{Mesh interconnect and operand injection}

Figure~\ref{fig:array} shows the interconnect. Internally, each PE reads the
\emph{registered} outputs of its four neighbours, so data advances by at most
one hop per step. External operands enter only at two edges: the north input
port feeds row~0 (one value per column) and the west input port feeds column~0
(one value per row). The south and east edges read a constant zero. All sixteen
PE registers are exposed for read-back.

\begin{figure*}[t]
  \centering
  % figures/array.tex — 4x4 PE mesh with edge operand injection (logical).
\begin{tikzpicture}[x=15mm, y=15mm]
  % PE grid: (row r, col c), r=0 top
  \foreach \r in {0,1,2,3} {
    \foreach \c in {0,1,2,3} {
      \node[pe] (p\r\c) at (\c, -\r) {\scriptsize$\r,\c$};
    }
  }
  % horizontal neighbour links (west->east)
  \foreach \r in {0,1,2,3} {
    \foreach \c/\cn in {0/1,1/2,2/3} {
      \draw[wire] (p\r\c) -- (p\r\cn);
    }
  }
  % vertical neighbour links (north->south)
  \foreach \c in {0,1,2,3} {
    \foreach \r/\rn in {0/1,1/2,2/3} {
      \draw[wire] (p\r\c) -- (p\rn\c);
    }
  }
  % north injection into row 0 (per column)
  \foreach \c in {0,1,2,3} {
    \draw[wire] (\c, 0.9) node[lbl, above]{$\mathrm{north}_\c$} -- (p0\c);
  }
  % west injection into column 0 (per row)
  \foreach \r in {0,1,2,3} {
    \draw[wire] (-0.9, -\r) node[lbl, left]{$\mathrm{west}_\r$} -- (p\r0);
  }
  % south / east edges read zero
  \foreach \c in {0,1,2,3} {
    \node[lbl, below] at (\c, -3.6) {$0$};
    \draw[wire, dashed] (p3\c) -- (\c, -3.55);
  }
  \foreach \r in {0,1,2,3} {
    \node[lbl, right] at (3.55, -\r) {$0$};
    \draw[wire, dashed] (p\r3) -- (3.5, -\r);
  }
\end{tikzpicture}

  \caption{The $4\times4$ mesh. Solid arrows are nearest-neighbour register
  connections; operands are injected only at the north edge (per column) and
  the west edge (per row), while the south and east edges read zero. This
  edge-only injection is the central constraint on dataflow mapping
  (Section~\ref{sec:dataflow}).}
  \label{fig:array}
\end{figure*}

This edge-only injection, combined with the single register per PE, is the
defining architectural constraint of the fabric: interior PEs receive fresh
host data only through their neighbours, so any computation that uses the whole
array must route operands inward from the edges.
